\newcommand{\auditoriaid}{A10 -- Compilación, formato y cierre}
% Preambulo comun de las auditorias del informe E1 (revision_e1/)
\documentclass[11pt,a4paper]{article}
\usepackage[utf8]{inputenc}
\usepackage[T1]{fontenc}
\usepackage[spanish,es-noshorthands]{babel}
\usepackage{lmodern}
\usepackage{microtype}
% Sin patrones de guionado en espanol en esta instalacion (ver A10): se
% desactiva el guionado automatico para evitar cortes con reglas del ingles.
\hyphenpenalty=10000
\exhyphenpenalty=10000
\usepackage[a4paper,margin=2.2cm]{geometry}
\usepackage{booktabs,array,longtable}
\usepackage{graphicx,caption,subcaption}
\usepackage{xcolor}
\usepackage{enumitem}
\usepackage{fancyhdr}
\usepackage{amsmath}
\usepackage{tikz}
\usetikzlibrary{arrows.meta,positioning}
\usepackage{natbib}
\usepackage{xurl}
\usepackage[hidelinks]{hyperref}
\urlstyle{tt}
\newcommand{\archivo}[1]{\url{#1}}
\newcolumntype{L}[1]{>{\raggedright\arraybackslash}p{#1}}
\definecolor{alta}{RGB}{180,30,30}
\definecolor{media}{RGB}{200,120,0}
\definecolor{baja}{RGB}{60,110,160}
\definecolor{cajafondo}{RGB}{245,247,250}
\newcommand{\sevA}{\textcolor{alta}{\textbf{Alta}}}
\newcommand{\sevM}{\textcolor{media}{\textbf{Media}}}
\newcommand{\sevB}{\textcolor{baja}{\textbf{Baja}}}
\setlength{\parskip}{0.35em}
\setlength{\emergencystretch}{3em}
\setlength{\parindent}{0pt}
\pagestyle{fancy}
\fancyhf{}
\fancyhead[L]{\small Auditoría del Informe del Entregable~1 -- \auditoriaid}
\fancyhead[R]{\small \thepage}
\renewcommand{\headrulewidth}{0.4pt}
% Entorno de hallazgos: ID | Severidad | Ubicacion | Hallazgo y evidencia | Correccion
\newenvironment{hallazgos}{%
\small
\begin{longtable}{@{}L{1.15cm}L{1.25cm}L{1.9cm}L{6.0cm}L{\dimexpr\textwidth-10.3cm-10.6\tabcolsep\relax}@{}}
\toprule
\textbf{ID} & \textbf{Sev.} & \textbf{Ubicación} & \textbf{Hallazgo y evidencia} & \textbf{Corrección aplicada} \\
\midrule\endfirsthead
\toprule
\textbf{ID} & \textbf{Sev.} & \textbf{Ubicación} & \textbf{Hallazgo y evidencia} & \textbf{Corrección aplicada} \\
\midrule\endhead
\bottomrule\endfoot
}{\end{longtable}}
% Macros compartidas con el informe revisado (las secciones propuestas las usan)
% Macros compartidas entre el informe revisado y las auditorias.
% Cifras verificadas contra los archivos del repositorio (ver A01/A04/A07):
\newcommand{\nUniverso}{102}   % source_id CMIP6 con tos/Omon en ESGF (00b, model_availability_priority.csv)
\newcommand{\nSeed}{41}        % publican los cuatro experimentos (config/models_seed_cmip6.csv)
\newcommand{\nSel}{40}         % completos + QC PASS (models_inventory_final.csv)
\newcommand{\nComb}{160}       % combinaciones modelo-experimento PASS (qc_report.csv)
\newcommand{\nNoEnc}{44}       % config/models_no_encontrado_44.csv
\newcommand{\nParcial}{18}     % 102 - 40 - 44
% Fecha de la portada: fija (no \today), editable en un solo lugar.
\newcommand{\fechaentrega}{28 de septiembre de 2026}
% Estado de codigo que documenta el E1 (antes de los cambios del E2)
\newcommand{\commitEone}{\texttt{8ee6184}}

% En las auditorias el texto propuesto se muestra fuera del informe: las
% referencias cruzadas se imprimen con el nombre de la seccion destino.
\makeatletter
\def\etq#1#2{\expandafter\def\csname etq@#1\endcsname{#2}}
\etq{sec:resumen}{Resumen}
\etq{sec:alcance}{Alcance}
\etq{sec:marco}{Revisión científica}
\etq{sec:antecedentes}{Antecedentes}
\etq{sec:referencias}{Literatura sobre TOE}
\etq{tab:referencias}{bibliografía TOE}
\etq{sec:similitudes}{Niño~3.4 frente a Niño~1+2}
\etq{sec:toe_metodos}{Elección de los métodos de TOE}
\etq{tab:metodos_toe}{métodos TOE}
\etq{eq:szabo_signal}{ecuación señal M1}
\etq{eq:szabo_V}{ecuación V(t)}
\etq{eq:szabo_T}{ecuación T(t)}
\etq{eq:szabo_toe}{ecuación TOE M1}
\etq{eq:tod}{ecuación ToD}
\etq{sec:datos}{Datos}
\etq{fig:cmip6_contribution}{contribución institucional}
\etq{sec:tabla_modelos}{Modelos seleccionados}
\etq{tab:modelos}{modelos}
\etq{sec:desviacion}{Selección de modelos}
\etq{sec:dominios_periodos}{Dominios y periodos}
\etq{fig:region_nino}{regiones Niño}
\etq{sec:arquitectura}{Arquitectura del pipeline}
\etq{fig:flujograma}{Flujograma}
\etq{sec:scripts}{Descripción de los pasos}
\etq{sec:paso00b}{Paso 00b}
\etq{sec:paso01}{Paso 01}
\etq{sec:paso02}{Paso 02}
\etq{sec:paso04}{Paso 04}
\etq{sec:paso05}{Paso 05}
\etq{sec:paso06}{Paso 06}
\etq{sec:paso07}{Paso 07}
\etq{sec:reporte}{Reporte de estado}
\etq{sec:resultados}{Resultados}
\etq{fig:qc}{matriz de estado}
\etq{tab:cuentas}{conteo de modelos}
\etq{fig:mapas2d}{mapas 2D}
\etq{fig:series}{series}
\etq{fig:boxplot}{boxplots}
\etq{sec:roadmap}{Próximos pasos}
\etq{tab:roadmap}{próximos pasos}
\etq{sec:conclusiones}{Conclusiones}
\etq{sec:recomendaciones}{Recomendaciones}
\makeatother

\makeatletter
\AtBeginDocument{\renewcommand{\ref}[1]{\@ifundefined{etq@#1}{\textit{[#1]}}{\textit{\csname etq@#1\endcsname}}}}
\makeatother
% Texto propuesto: secciones degradadas un nivel y en caja
\newenvironment{propuesto}{%
  \par\medskip\noindent\textcolor{baja}{\rule{\textwidth}{0.8pt}}\par
  \noindent{\small\textbf{Inicio del texto propuesto para el informe revisado}}\par\medskip
  \begingroup
  \let\section\subsection\let\subsection\subsubsection\let\subsubsection\paragraph
}{%
  \endgroup
  \par\noindent{\small\textbf{Fin del texto propuesto}}\par
  \noindent\textcolor{baja}{\rule{\textwidth}{0.8pt}}\par\medskip}
% Recuadro de actualizacion posterior (decisiones del autor)
\newcommand{\actualizacion}[1]{\par\medskip\noindent\fcolorbox{media}{media!8}{\parbox{\dimexpr\linewidth-2\fboxsep-2\fboxrule\relax}{\small\textbf{Actualización del 29-09-2026 (decisiones del autor).} #1}}\par\medskip}

\begin{document}

\begin{center}
{\Large\bfseries Auditoría A10: Compilación, formato y cierre del Entregable~1}\\[0.4em]
{\normalsize \texttt{informe\_e1\_smnv3.tex} (commit \texttt{330fd2e}) y estado del repositorio}\\[0.2em]
{\small Revisión: \fechaentrega}
\end{center}

\section{Objeto}

Se audita cómo se compila y cómo se ve el informe, y qué falta en el repositorio para que el Entregable~1 quede cerrado y reproducible. El informe original se recompiló en una copia aparte (\archivo{revision_e1/trabajo/compila_original/}) con \texttt{pdflatex}, \texttt{bibtex} y dos pasadas más, y se inspeccionaron el \emph{log}, las fuentes incrustadas (\texttt{pdffonts}) y las páginas renderizadas.

\section{Hallazgos de compilación y formato}

\begin{hallazgos}
A10-1 & \sevA & Fuentes & El PDF incrusta 18 fuentes Type~3 (mapa de bits) y solo 6 Type~1. Con \texttt{fontenc} T1 y sin \texttt{lmodern} ni \texttt{cm-super}, el texto se rasteriza: se ve borroso al ampliar e imprime peor. & \texttt{\textbackslash usepackage\{lmodern\}}: todas las fuentes pasan a ser vectoriales. \\
A10-2 & \sevA & Guionado & \emph{Log}: ``No hyphenation patterns were preloaded for the language `Spanish' \ldots\ I will use the patterns preloaded for \textbackslash language=0''. En esta máquina no está instalado \texttt{texlive-lang-spanish}, así que las palabras se cortan con reglas del inglés. & Informe revisado: guionado automático desactivado y compensado con \texttt{microtype} y \texttt{\textbackslash emergencystretch}. Recomendado: \texttt{sudo apt install texlive-lang-spanish}, y luego volver a activar el guionado. \\
A10-3 & \sevM & Pasadas & Con \texttt{pdflatex}, \texttt{bibtex}, \texttt{pdflatex}, \texttt{pdflatex} queda ``Label(s) may have changed''; hace falta una cuarta pasada (\texttt{longtable} + \texttt{hyperref}). & Script de compilación con cuatro pasadas de \texttt{pdflatex}. \\
A10-4 & \sevM & \texttt{hyperref} & ``destination with the same identifier (name\{page.1\})'': la portada y el índice son ambos página~1. En consecuencia, el encabezado numera ``11'' la página~12 del PDF. & \texttt{\textbackslash pagenumbering}: letras para la portada, romanos para los índices y arábigos desde el Resumen. \\
A10-5 & \sevM & Cuadro de modelos & 6 \emph{overfull hbox} en el encabezado del \texttt{longtable} (``Código'', ``Realización'', ``Resolución'' en columnas de 1,1 y 1,9\,cm). & Anchos recalculados y cuadro en \texttt{\textbackslash small}. \\
A10-6 & \sevM & Figuras & \archivo{informe/figuras/} enlaza con enlaces simbólicos a \archivo{figures/}, que no se versiona y se regenera en cada corrida. El PDF cambia según el estado de la máquina que compila (A04-5, A07-1). & Informe revisado: figuras como archivos reales en \archivo{revision_e1/figuras/}. \\
A10-7 & \sevM & Fecha & \verb|\today| en la portada: cada recompilación cambia la fecha del entregable. & Macro \verb|\fechaentrega| (a confirmar). \\
A10-8 & \sevB & Codificación & Mezcla de acentos escritos como \verb|\'a| y en UTF-8 directo; \texttt{pdftitle} sin tildes (``Informe tecnico \ldots\ Nino'') junto a partes con tildes. & UTF-8 en todo el texto; título del PDF con tildes. \\
A10-9 & \sevB & Código fuente & Quedan 5 comentarios \texttt{\% REVISAR} y código TikZ comentado dentro del documento entregable. & Resueltos (A01, A02, A03, A05, A08) y eliminados. \\
A10-10 & \sevB & Anexo del informe & El árbol del repositorio del anexo omite archivos existentes (\texttt{pipeline\_config.py}, \texttt{preflight\_report.py}, \texttt{pipeline\_timing.py}, \archivo{bundle_deliverable_update.py}, \archivo{models_copernicus_ssp245_whitelist.csv}) y remite a una ``Sección~8'' que no corresponde. & Se reemplaza por el documento de anexos separado, con el árbol actualizado. \\
A10-11 & \sevB & Unidades & El símbolo de grado aparece como \texttt{\textbackslash textdegree} en unas partes y como ``°'' en otras. & Uniforme (°). \\
\end{hallazgos}

\section{Acciones pendientes en el repositorio para cerrar el E1}

Ninguna de estas acciones se ejecutó en esta revisión. Todas modifican el repositorio y requieren la decisión del autor; no se hizo ningún \emph{commit} ni \emph{push}.

\begin{enumerate}[leftmargin=1.5em]
\item \textbf{Etiquetar el estado del E1} (opcional): el informe ya no cita commits y describe la ventana vigente. Una etiqueta sobre el commit que se entregue sigue siendo útil como referencia interna.
\item \textbf{Congelar las figuras}: reemplazar los enlaces simbólicos de \archivo{informe/figuras/} por copias reales de \archivo{revision_e1/figuras/} (ventana vigente, etiquetas en español, mapa de regiones global).
\item \textbf{Corregir los scripts de gráficos}: redondeo del año decimal en \texttt{plot\_boxplot\_comparison.py} y \texttt{plot\_box\_series.py} (\texttt{np.round} → \texttt{np.floor}) y etiquetas en español (A07). El cambio está probado en \archivo{revision_e1/trabajo/e1_snapshot/}.
\item \textbf{Revisar la idempotencia de las figuras}: un PNG existente no se regenera aunque cambie el inventario (A07-1). Conviene regenerar cuando el inventario es más reciente que la figura.
\item \textbf{Actualizar comentarios y archivos desactualizados}: cabecera de \texttt{run.sh} (A05-7); \texttt{config/models\_missing\_from\_esgf.csv}, en el que 3 de sus 17 modelos ya están en el \emph{seed} (A04-2).
\item \textbf{Bibliografía}: adoptar \archivo{revision_e1/informe/references_rev.bib} (A09).
\item \textbf{Entorno de compilación}: instalar \texttt{texlive-lang-spanish} y agregar \texttt{lmodern} al informe original si se sigue usando.
\item \textbf{Limpieza}: \archivo{informe/borradores/informe_e1_smnv3.pdf} (no versionado, versión del 22-09) puede borrarse por su nombre exacto o renombrarse como borrador.
\end{enumerate}

\section{Estructura del informe revisado}

\begin{center}\small
\begin{tabular}{@{}lL{7.5cm}r@{}}
\toprule
Sección & Contenido & Auditoría \\
\midrule
Portada, índices & Mismo formato de tapa; fecha fija; numeración corregida & A01, A10 \\
1--2 & Resumen y Alcance & A01 \\
3 & Revisión científica (antecedentes, literatura, implicancias, métodos TOE) & A02, A03 \\
4 & Datos (fuentes, modelos, selección, dominios y periodos) & A04 \\
5 & Arquitectura del \emph{pipeline} (el uso pasa al anexo) & A05 \\
6 & Descripción de los pasos 00--07 & A06 \\
7 & Resultados & A07 \\
8--10 & Próximos pasos, conclusiones y recomendaciones & A08 \\
Bibliografía & \archivo{references_rev.bib} & A09 \\
\bottomrule
\end{tabular}
\end{center}

El documento de anexos (uso del \emph{pipeline}) se entrega aparte, como pidió el autor.

\end{document}
